\documentclass[10pt,twoside]{article}


\usepackage{iftex}

\ifPDFTeX
\usepackage[utf8]{inputenc}
\usepackage[T5]{fontenc}
\usepackage[vietnamese]{babel}
\usepackage{mathptmx}
\usepackage[scaled=0.92]{helvet}
\newcommand{\arialfont}{\fontfamily{phv}\selectfont}
\else
\usepackage[vietnamese]{babel}
\usepackage{fontspec}
\defaultfontfeatures{Ligatures=TeX}

\IfFontExistsTF{Times New Roman}
{\setmainfont{Times New Roman}}
{\setmainfont{TeX Gyre Termes}}

\IfFontExistsTF{Arial}
{\setsansfont{Arial}\newfontfamily\arialfont{Arial}}
{\setsansfont{TeX Gyre Heros}\newfontfamily\arialfont{TeX Gyre Heros}}
\fi

% ------------------ Page layout theo template FAIR ------------------
\usepackage[
paperwidth=20.5cm,
paperheight=29.5cm,
top=2.2cm,
bottom=1.8cm,
left=1.8cm,
right=1.8cm,
headheight=23pt,
headsep=0.55cm,
footskip=1cm
]{geometry}

% ------------------ Packages ------------------
\usepackage{amsmath,amssymb}
\usepackage{graphicx}
\usepackage{array}
\usepackage{booktabs}
\usepackage{caption}
\usepackage{subcaption}
\usepackage{float}
\usepackage{enumitem}
\usepackage{fancyhdr}
\usepackage{titlesec}
\usepackage{indentfirst}
\usepackage{etoolbox}
\usepackage{url}

% ------------------ Paragraph style ------------------
\setlength{\parindent}{1cm}
\setlength{\parskip}{6pt}
\linespread{1.0}
\emergencystretch=2em
\raggedbottom

% ------------------ Paper metadata ------------------
\newcommand{\PaperTitle}{XẾP HẠNG TÍN DỤNG DOANH NGHIỆP DỰA TRÊN MÔ HÌNH MÁY HỌC}
\newcommand{\PaperShortTitle}{XẾP HẠNG TÍN DỤNG DOANH NGHIỆP DỰA TRÊN MÔ HÌNH MÁY HỌC}
\newcommand{\PaperAuthors}{Lê Nguyễn Thành Tài\textsuperscript{1}}
\newcommand{\PaperShortAuthors}{Lê Nguyễn Thành Tài}
\newcommand{\PaperAffiliations}{%
\textsuperscript{1}Khoa Công nghệ Thông tin, Trường Đại học Sư phạm kỹ thuật Vĩnh Long
}
\newcommand{\PaperEmail}{22022001@st.vlute.edu.vn}
\newcommand{\FairConference}{}

% ------------------ Header/Footer ------------------
\pagestyle{fancy}
\fancyhf{}

% Trang chẵn: số trang bên trái, tên bài bên phải.
% Trang lẻ: tác giả bên trái, số trang bên phải.
\fancyhead[LE]{\fontsize{9pt}{11pt}\selectfont\thepage}
\fancyhead[RE]{\fontsize{9pt}{11pt}\selectfont\PaperShortTitle}
\fancyhead[LO]{\fontsize{9pt}{11pt}\selectfont\PaperShortAuthors}
\fancyhead[RO]{\fontsize{9pt}{11pt}\selectfont\thepage}

% Template có đường kẻ mảnh dưới header ở các trang sau.
\renewcommand{\headrulewidth}{0.4pt}
\renewcommand{\footrulewidth}{0pt}

\fancypagestyle{firstpage}{%
\fancyhf{}%
\fancyhead[C]{\fontsize{9pt}{11pt}\selectfont\FairConference}%
\renewcommand{\headrulewidth}{0.4pt}
\renewcommand{\footrulewidth}{0pt}%
}

% ------------------ Title block ------------------
\newcommand{\makefairtitle}{%
\thispagestyle{firstpage}%
\begin{center}
{\arialfont\bfseries\fontsize{14pt}{16pt}\selectfont\MakeUppercase{\PaperTitle}\par}
\vspace{3pt}
{\bfseries\normalsize\PaperAuthors\par}
\vspace{3pt}
{\fontsize{9pt}{11pt}\selectfont\PaperAffiliations\par}
\vspace{3pt}
{\itshape\fontsize{9pt}{11pt}\selectfont\PaperEmail\par}
\end{center}
\vspace{6pt}%
}

% ------------------ Abstract / Keywords ------------------
\newcommand{\FairDash}{\unskip\textemdash\space}

\newenvironment{fairabstract}{%
\begingroup
\fontsize{9pt}{11pt}\selectfont
\itshape
\setlength{\parindent}{1cm}%
\setlength{\parskip}{6pt}%
}{%
\par\endgroup
}

\newcommand{\fairabstracttext}[1]{%
\begin{fairabstract}
\textbf{TÓM TẮT}\FairDash #1
\end{fairabstract}
}

\newcommand{\fairkeywords}[1]{%
\begin{fairabstract}
\textbf{Từ khóa}\FairDash #1
\end{fairabstract}
}

% ------------------ Section numbering ------------------
\setcounter{secnumdepth}{4}

\renewcommand{\thesection}{\Roman{section}}
\renewcommand{\thesubsection}{\Alph{subsection}}
\renewcommand{\thesubsubsection}{\arabic{subsubsection}}
\renewcommand{\theparagraph}{\alph{paragraph}}

\titleformat{\section}
{\centering\normalsize\bfseries}
{\thesection.}{0.5em}{}

\titleformat{\subsection}
{\normalsize\bfseries\itshape}
{\thesubsection.}{0.5em}{}

\titleformat{\subsubsection}
{\normalsize\bfseries\itshape}
{\thesubsubsection.}{0.5em}{}

\titleformat{\paragraph}[block]
{\normalsize}
{\theparagraph)}{0.5em}{}

\titlespacing*{\section}{0pt}{6pt}{6pt}
\titlespacing*{\subsection}{0pt}{6pt}{3pt}
\titlespacing*{\subsubsection}{0pt}{6pt}{3pt}
\titlespacing*{\paragraph}{0pt}{6pt}{3pt}

% ------------------ Figure/Table captions ------------------
\renewcommand{\figurename}{Hình}
\renewcommand{\tablename}{Bảng}

\captionsetup{
font=small,
labelsep=period,
justification=centering,
singlelinecheck=true,
skip=3pt
}


% ------------------ Lists ------------------
\setlist[itemize]{
leftmargin=1.27cm,
itemsep=0pt,
topsep=0pt,
parsep=0pt
}

\setlist[enumerate]{
leftmargin=1.27cm,
itemsep=0pt,
topsep=0pt,
parsep=0pt
}

% ------------------ Bibliography ------------------
% Giữ tiêu đề TÀI LIỆU THAM KHẢO theo kiểu section nhưng không đánh số.
\makeatletter
\renewenvironment{thebibliography}[1]{%
\section*{TÀI LIỆU THAM KHẢO}%
\fontsize{9pt}{11pt}\selectfont
\hbadness=10000
\setlength{\parskip}{0pt}%
\list{[\arabic{enumiv}]}{%
\settowidth\labelwidth{[#1]}%
\leftmargin\labelwidth
\advance\leftmargin\labelsep
\usecounter{enumiv}%
\let\p@enumiv\@empty
\renewcommand{\theenumiv}{\arabic{enumiv}}%
}%
\sloppy
\clubpenalty4000
\widowpenalty4000
}{%
\endlist
}
\makeatother

\urlstyle{same}

% =========================================================
% DOCUMENT
% =========================================================
\begin{document}

\makefairtitle

\fairabstracttext{%
Xếp hạng tín dụng doanh nghiệp có vai trò quan trọng trong đánh giá sức khỏe tài chính, khả năng thực hiện nghĩa vụ nợ và cảnh báo sớm rủi ro. Nghiên cứu này xây dựng một hệ thống xếp hạng tín dụng dựa trên dữ liệu tài chính và so sánh các mô hình LightGBM, XGBoost, LSTM, TCN, PatchTST, Transformer-LSTM và GraphSAGE. Tập dữ liệu thực nghiệm gồm 8.680 quan sát của 1.623 doanh nghiệp trong giai đoạn 2005--2016; mỗi quan sát được biểu diễn bằng 12 chỉ số thuộc các nhóm thanh khoản, đòn bẩy, sinh lời, hiệu quả hoạt động và dòng tiền, đồng thời được phân loại thành ba nhóm Distressed, High Yield và Investment Grade. Bên cạnh các mô hình riêng lẻ, nghiên cứu triển khai Fuzzy Choquet Integral, Gompertz Fuzzy Ranking và đề xuất Dynamic Model Fusion kết hợp Transformer-LSTM với GraphSAGE. Khi hai mô hình thành phần bất đồng, cơ chế Dynamic Classifier Selection lựa chọn mô hình phù hợp theo năng lực trên tập xác thực và độ tin cậy đối với từng mẫu. Kết quả cho thấy Dynamic Model Fusion đạt hiệu suất tổng thể cao nhất, với Accuracy 92,98\%, Precision 92,67\%, Recall 92,98\% và F1-Score 92,60\%. GradientSHAP và LIME được tích hợp để giải thích ảnh hưởng của các đặc trưng tài chính; kết quả nhấn mạnh vai trò của biên lợi nhuận, tỷ lệ nợ trên vốn chủ sở hữu và khả năng thanh khoản. Hệ thống nhờ đó vừa duy trì hiệu quả phân loại, vừa tăng tính minh bạch và khả năng hỗ trợ ra quyết định trong đánh giá rủi ro tín dụng doanh nghiệp.
}

\fairkeywords{%
Xếp hạng tín dụng doanh nghiệp, Máy học, Học sâu, Ensemble Learning, Explainable AI.
}

\section{GIỚI THIỆU}
\subsection{Giới thiệu bài toán}
Xếp hạng tín dụng doanh nghiệp là một trong những bài toán quan trọng trong lĩnh vực tài chính, phản ánh khả năng thanh toán, mức độ rủi ro và sức khỏe tài chính của doanh nghiệp. Trong bối cảnh kinh tế toàn cầu tiếp tục chịu tác động từ biến động lãi suất và thị trường vốn, rủi ro vỡ nợ doanh nghiệp đang trở thành thách thức lớn đối với hệ thống tài chính. Theo S\&P Global Ratings, năm 2025 toàn cầu ghi nhận 117 vụ vỡ nợ, tuy số lượng giảm nhưng tổng giá trị nợ mất khả năng thanh toán vẫn vượt 140 tỷ USD \cite{noauthor_default_nodate}. Tại khu vực Đông Á mới nổi, quy mô thị trường trái phiếu bằng đồng nội tệ đạt khoảng 29,5 nghìn tỷ USD vào cuối năm 2025, cho thấy nhu cầu đánh giá chất lượng tín dụng ngày càng cao \cite{asian_development_bank_asia_2017}. Riêng tại Việt Nam, thị trường trái phiếu doanh nghiệp phục hồi đạt khoảng 1,42 triệu tỷ đồng, tương đương gần 11\% GDP, trong đó lượng phát hành mới năm 2025 tăng 32\% và đạt hơn 624 nghìn tỷ đồng \cite{noauthor_tong_nodate}. Những số liệu này cho thấy việc xây dựng một hệ thống xếp hạng tín dụng minh bạch và đáng tin cậy là yêu cầu cấp thiết. Tuy nhiên, việc đánh giá rủi ro tín dụng doanh nghiệp trong thực tế vẫn gặp nhiều khó khăn do dữ liệu tài chính có tính biến động cao, quan hệ giữa các chỉ số tài chính thường phức tạp và phi tuyến. Các phương pháp đánh giá truyền thống phụ thuộc nhiều vào kinh nghiệm chuyên gia, trong khi các mô hình học máy hiện đại dù có khả năng dự báo tốt nhưng thường bị xem là “hộp đen”, khó giải thích nguyên nhân dẫn đến kết quả xếp hạng. Điều này làm giảm mức độ tin cậy khi áp dụng trong các quyết định tài chính có yêu cầu cao về tính minh bạch. Do đó, bài toán đặt ra là xây dựng một hệ thống xếp hạng tín dụng doanh nghiệp dựa trên các mô hình máy học và học sâu, có khả năng phân loại mức độ rủi ro chính xác, đồng thời giải thích được vai trò của từng chỉ số tài chính đối với kết quả dự báo.
\subsection{Những nghiên cứu liên quan}
Trương Thị Thùy Dương và Lê Hải Trung \cite{thuy_ung_2023} đã ứng dụng các mô hình Logistic Regression, Decision Tree, Random Forest, K-Nearest Neighbor, Naïve Bayes và XGBoost để dự báo rủi ro phá sản của 300 doanh nghiệp Việt Nam giai đoạn 2017–2019. Kết quả cho thấy XGBoost đạt hiệu quả cao nhất với độ chính xác 87.73\%, tiếp theo là Random Forest với 83.49\%. Kofune và cộng sự \cite{kofune_enhancing_2026} đề xuất mô hình Transformer-LSTM cho phân loại xếp hạng tín dụng doanh nghiệp, kết hợp khả năng khai thác phụ thuộc dài hạn của Transformer và xử lý chuỗi thời gian của LSTM, đạt độ chính xác 89.96\%. Feng và cộng sự \cite{feng_every_2020} đề xuất CCR-CNN bằng cách chuyển đổi dữ liệu tài chính thành dạng ảnh và sử dụng CNN để khai thác tương tác đặc trưng, đạt độ chính xác 95\% trên bộ dữ liệu doanh nghiệp niêm yết tại Trung Quốc. Tiếp đó, Feng và cộng sự \cite{feng_corporate_2024} phát triển mô hình HHGNN nhằm khai thác cấu trúc phân cấp, tính không đồng nhất và dữ liệu chưa gán nhãn trong mạng lưới doanh nghiệp; mô hình đạt độ chính xác 97.557\%, vượt trội so với LR, SVM, MLP, XGBoost, HGT, CCR-CNN, CCR-GNN và ASSL4CCR. Chen và Long \cite{chen_novel_2020} xây dựng mô hình SMAGRU dựa trên cơ chế Self-Attention cho xếp hạng tín dụng doanh nghiệp end-to-end và đạt độ chính xác 93.33\%. Morthens \cite{morthens_credit_nodate} đề xuất mô hình chuỗi thời gian đa phương thức kết hợp RNN và GNN để khai thác đồng thời dữ liệu tài chính theo thời gian và quan hệ doanh nghiệp; trong đó GRU đạt 75.25\% và GRU-GCNGRU đạt 74.44\%. Liu và Kok \cite{liu_prediction_2025} đề xuất mô hình Heterogeneous Topological Graph Neural Networks (HTGNN) cho dự báo xếp hạng tín dụng ngân hàng, đạt độ chính xác 70.57\%. Andersson \cite{andersson_masters_nodate} xây dựng hệ thống học máy dự báo xếp hạng tín dụng doanh nghiệp trên bộ dữ liệu khoảng 3 triệu dòng với hơn 40 biến tài chính; kết quả cho thấy Gradient Boosting đạt hiệu quả tốt nhất, với độ chính xác xấp xỉ 90\% đối với các tổ chức xếp hạng lớn như Standard \& Poor's, Moody's và Fitch.Fang và cộng sự \cite{fang_design_2025} đã đề xuất hệ thống dự báo rủi ro tín dụng đa phương thức dựa trên phân tích văn bản tài chính. Nghiên cứu kết hợp oversampling để xử lý mất cân bằng dữ liệu, chọn lọc đặc trưng và mô hình lai CNN-BiLSTM-AM nhằm khai thác thông tin ngữ nghĩa và cảm xúc quản lý. Kết quả cho thấy phương pháp đề xuất đạt hiệu quả dự báo tốt hơn các mô hình phân tích văn bản truyền thống đạt độ chính xác 88\%.
\section{CÁC CÔNG VIỆC LIÊN QUAN}

\subsection{Đặc điểm của Xếp hạng tín dụng doanh nghiệp}

Xếp hạng tín dụng doanh nghiệp là quá trình đánh giá mức độ tin cậy tài chính và khả năng thực hiện nghĩa vụ nợ của doanh nghiệp trong một khoảng thời gian nhất định \cite{phu_nghi_nodate}. Kết quả xếp hạng phản ánh sức khỏe tài chính, khả năng thanh toán và mức độ rủi ro vỡ nợ, qua đó hỗ trợ ngân hàng, nhà đầu tư và tổ chức quản lý trong quyết định cấp tín dụng, đầu tư trái phiếu và giám sát rủi ro tài chính \cite{hai_mo_2023}. Thông thường, doanh nghiệp có khả năng sinh lời tốt, dòng tiền ổn định và đòn bẩy tài chính thấp sẽ được xếp vào nhóm tín nhiệm cao; ngược lại, doanh nghiệp có dòng tiền yếu, nợ vay lớn hoặc nguy cơ mất khả năng thanh toán sẽ thuộc nhóm rủi ro cao \cite{thuy_ung_2023}. Trong nghiên cứu này, xếp hạng tín dụng doanh nghiệp được phân thành ba nhóm gồm Investment Grade (IG), High Yield (HY) và Distressed. Nhóm IG đại diện cho doanh nghiệp có rủi ro thấp và khả năng thanh toán ổn định; nhóm HY thể hiện doanh nghiệp có tiềm năng sinh lợi nhưng đi kèm rủi ro cao hơn; trong khi nhóm Distressed biểu thị doanh nghiệp có tình trạng tài chính yếu và nguy cơ vỡ nợ lớn \cite{hung_xep_2023}. Bài toán xếp hạng tín dụng chủ yếu dựa trên các chỉ số tài chính về thanh khoản, đòn bẩy, khả năng sinh lời và dòng tiền, nhưng mối quan hệ giữa các chỉ số này với rủi ro tín dụng thường phi tuyến và biến động theo thời gian \cite{hai_mo_2023}. Do đó, việc ứng dụng các mô hình máy học, học sâu và Explainable AI là cần thiết nhằm nâng cao độ chính xác dự báo, đồng thời giải thích rõ vai trò của từng chỉ số tài chính trong kết quả xếp hạng \cite{gramegna_shap_2021}.
\subsection{Các mô hình mạng dùng trong huấn luyện}

\subsubsection{Mô hình LightGBM (Light Gradient Boosting Machine)}
Mô hình LightGBM (Light Gradient Boosting Machine) \cite{ke_lightgbm_2017} được Guolin Ke và các cộng sự tại Microsoft Research giới thiệu vào năm 2017. Tương tự như XGBoost, LightGBM là một thuật toán học tập hợp (Ensemble Learning) mạnh mẽ thuộc họ Cây quyết định tăng cường theo gradient (GBDT). Tuy nhiên, LightGBM được thiết kế với mục tiêu cốt lõi là giải quyết triệt để bài toán về hiệu suất tính toán và giới hạn bộ nhớ khi phải xử lý các tập dữ liệu có kích thước khổng lồ (Big Data) và không gian đặc trưng nhiều chiều.
\subsubsection{Mô hình XGBoost (eXtreme Gradient Boosting)}
Mô hình XGBoost (eXtreme Gradient Boosting) \cite{chen_xgboost_2016} được Tianqi Chen, Carlos Guestrin và các cộng sự tại Đại học Washington giới thiệu lần đầu tại hội nghị KDD 2016 thông qua bài báo khoa học “XGBoost: A Scalable Tree Boosting System”. Đây là một thuật toán học tập hợp (Ensemble Learning) tiên tiến, được phát triển dựa trên nền tảng của Cây quyết định tăng cường theo gradient (GBDT), nhưng được tối ưu hóa vượt trội về tài nguyên tính toán và khả năng kiểm soát độ phức tạp.
\subsubsection{Mạng tích chập thời gian (Temporal Convolutional Network - TCN)}
Mô hình Mạng tích chập thời gian (TCN) \cite{bai_empirical_2018} được đề xuất trong bài báo nghiên cứu "An Empirical Evaluation of Generic Convolutional and Recurrent Networks for Sequence Modeling" của Shaojie Bai, J. Zico Kolter và Vladlen Kol-tun. TCN là một kiến trúc tích chập thuần túy được thiết kế riêng cho các tác vụ xử lý dữ liệu tuần tự, lấy cảm hứng từ việc tối ưu hóa khả năng tính toán song song của CNN để thay thế cho các kiến trúc hồi tiếp truyền thống như RNN hay LSTM.
\subsubsection{Mạng nơ-ron bộ nhớ ngắn hạn dài (Long Short-Term Memory - LSTM)}
Mô hình Mạng nơ-ron bộ nhớ ngắn hạn dài (LSTM) \cite{hochreiter_long_1997} là một biến thể kiến trúc bậc cao của Mạng nơ-ron hồi tiếp (RNN), được Sepp Hochreiter và Jür-gen Schmidhuber giới thiệu lần đầu vào năm 1997. Mục tiêu cốt lõi của LSTM là khắc phục triệt để bài toán "triệt tiêu đạo hàm" (vanishing gradient) – một nhược điểm chí mạng khiến các mạng RNN truyền thống không thể học được từ các chuỗi dữ liệu dài.
\subsubsection{Kiến trúc PatchTST (Patch Time Series Transformer)}
Kiến trúc PatchTST \cite{nie_time_2023} được đề xuất trong bài báo nghiên cứu "A Time Series is Worth 64 Words: Long-term Forecasting with Transformers" của Yuqi Nie, Nam H. Nguyen và các cộng sự tại IBM Research, được công bố tại ICLR 2023. PatchTST là một mô hình tiên phong lấy cảm hứng từ sự thành công của Vi-sion Transformers (ViT), nhưng được tối ưu hóa đặc biệt để vượt qua các giới hạn về chi phí tính toán và khả năng nắm bắt ngữ nghĩa cục bộ trong dữ liệu chuỗi thời gian.
\subsubsection{Mô hình lai T-LSTM(Transformer-Long Short-Term Memory)}
Mô hình lai T-LSTM (Transformer-Long Short-Term Memory) \cite{kofune_enhancing_2026} là kiến trúc kết hợp giữa Transformer và LSTM nhằm khai thác hiệu quả dữ liệu tài chính dạng chuỗi thời gian. Trong bài toán xếp hạng tín dụng doanh nghiệp, dữ liệu tài chính không chỉ phản ánh trạng thái tại một thời điểm mà còn thể hiện xu hướng biến động của doanh nghiệp qua nhiều kỳ liên tiếp. Vì vậy, mô hình T-LSTM được sử dụng để học đồng thời các quan hệ dài hạn và sự phụ thuộc tuần tự trong chuỗi chỉ số tài chính.
\subsubsection{Mạng nơ ron đồ thị GraphSAGE (Graph Sample and Aggregate)}
Mô hình đồ thị được sử dụng trong nghiên cứu này dựa trên ý tưởng từ công trình Inductive Representation Learning on Large Graphs của Hamilton, Ying và Leskovec \cite{hamilton_inductive_2018}. Trong nghiên cứu này, nhóm tác giả đề xuất GraphSAGE một khung học biểu diễn nút theo hướng quy nạp, có khả năng tạo embedding cho các nút mới bằng cách học hàm tổng hợp thông tin từ đặc trưng của nút và các nút lân cận, thay vì học trực tiếp một vector embedding cố định cho từng nút như các phương pháp transductive truyền thống.
\subsubsection{Độ đo đánh giá mô hình} 
\paragraph{Accuracy} 
Accuracy \cite{grandini_metrics_2020} hay độ chính xác, là tỷ lệ giữa số mẫu được mô hình dự đoán đúng trên tổng số mẫu trong tập dữ liệu. Đây là một độ đo đơn giản, dễ hiểu và thường được sử dụng trong các bài toán phân loại để đánh giá hiệu suất tổng quát của mô hình. Công thức tính Accuracy được biểu diễn như sau: 
\begin{equation} Accuracy = \frac{TP + TN}{TP + TN + FP + FN} 
\end{equation} Trong đó: 
\begin{itemize} 
\item $TP$ (True Positive): số lượng mẫu dương tính được dự đoán đúng. 
\item $TN$ (True Negative): số lượng mẫu âm tính được dự đoán đúng. 
\item $FP$ (False Positive): số lượng mẫu âm tính nhưng bị dự đoán nhầm là dương tính. 
\item $FN$ (False Negative): số lượng mẫu dương tính nhưng bị dự đoán nhầm là âm tính. 
\end{itemize}
\paragraph{Precision} 
Precision \cite{grandini_metrics_2020} đo lường tỷ lệ dự đoán dương chính xác trên tổng số dự đoán dương. Điều này đặc biệt quan trọng trong các ứng dụng mà dự đoán dương giả (false positives) có thể gây hậu quả nghiêm trọng.
\begin{equation} 
\mathrm{Precision}=\frac{TP}{TP+FP}
\end{equation}
\paragraph{Recall} 
Recall \cite{grandini_metrics_2020} đo lường tỷ lệ dự đoán dương chính xác trên tổng số các mẫu thực tế là dương. Recall quan trọng trong các tình huống cần phát hiện tất cả các trường hợp dương, ngay cả khi điều đó dẫn đến việc dự đoán một số trường hợp âm là dương (false negatives).
\begin{equation} 
\mathrm{Recall}=\frac{TP}{TP+FN}
\end{equation}
\paragraph{F1-Score} 
F1-Score \cite{grandini_metrics_2020} là trung bình điều hòa của Precision và Recall. Nó cung cấp một cách đánh giá cân bằng giữa độ chính xác của các dự đoán dương và khả năng phát hiện đầy đủ các mẫu dương.
\begin{equation} 
F1=2\times\frac{\mathrm{Precision}\times\mathrm{Recall}}{\mathrm{Precision}+\mathrm{Recall}}
\end{equation}
\paragraph{AUC-ROC} 
ROC Curve biểu diễn mối quan hệ giữa True Positive Rate (Recall) và False Positive Rate. AUC là diện tích dưới đường cong ROC và thể hiện khả năng phân biệt giữa các lớp của mô hình \cite{fawcett_introduction_2006}. Một mô hình với AUC gần 1 là một mô hình tốt, trong khi AUC gần 0.5 cho thấy mô hình hoạt động như ngẫu nhiên.
\begin{equation} 
\mathrm{TPR}=\frac{TP}{TP+FN}\quad;\quad \mathrm{FPR}=\frac{FP}{FP+TN}
\end{equation}

\subsection{Các Phương pháp học kết hợp}
\subsubsection{Fuzzy Choquet Integral} Fuzzy Choquet Integral là phương pháp học kết hợp dựa trên độ đo mờ, được sử dụng để tổng hợp kết quả dự đoán từ nhiều mô hình thành phần. Khác với trung bình có trọng số, phương pháp này không chỉ xét mức độ quan trọng riêng lẻ của từng mô hình mà còn biểu diễn được quan hệ tương tác giữa các mô hình trong ensemble \cite{zadeh_fuzzy_1965, choquet_theory_1954}. Giả sử $x_1, x_2, \ldots, x_n$ là xác suất dự đoán của $n$ mô hình, được sắp xếp tăng dần $x_{(1)} \leq x_{(2)} \leq \cdots \leq x_{(n)}$. Tích phân Choquet rời rạc được tính như sau: \begin{equation} C_{\mu}(x) = \sum_{i=1}^{n} \left(x_{(i)} - x_{(i-1)}\right)\mu(A_{(i)}), \end{equation} trong đó $x_{(0)}=0$, $A_{(i)}=\{(i),(i+1),\ldots,(n)\}$ và $\mu(A_{(i)})$ là độ đo mờ biểu diễn mức độ quan trọng của tập các mô hình tương ứng.

\subsubsection{Gompertz Fuzzy Ranking} 
Gompertz Fuzzy Ranking là phương pháp học kết hợp sử dụng hàm Gompertz để gán trọng số mờ cho các mô hình dựa trên thứ hạng xác suất dự báo. Thay vì trung bình trực tiếp đầu ra của các mô hình, phương pháp này ưu tiên các mô hình có dự đoán tin cậy hơn và giảm ảnh hưởng của các mô hình có xác suất thấp \cite{kundu_fuzzy_2021}. Giả sử $p_i^{(c)}$ là xác suất dự báo của mô hình thứ $i$ đối với lớp $c$. Các xác suất được sắp xếp theo thứ hạng $r_i$, trong đó mô hình có xác suất cao hơn sẽ có thứ hạng tốt hơn. Trọng số mờ Gompertz được xác định như sau: 
\begin{equation} 
w_i = \frac{\exp \left(-\alpha e^{-\beta r_i}\right)} {\sum_{j=1}^{M} \exp \left(-\alpha e^{-\beta r_j}\right)} 
\end{equation} 
trong đó $M$ là số mô hình thành phần, $\alpha$ và $\beta$ là các tham số điều chỉnh độ dốc của hàm Gompertz. Điểm dự báo tổng hợp cho lớp $c$ được tính bởi: 
\begin{equation} 
S_c = \sum_{i=1}^{M} w_i p_i^{(c)} 
\end{equation} 
Nhãn dự đoán cuối cùng được xác định theo lớp có điểm tổng hợp lớn nhất. Trong bài toán xếp hạng tín dụng doanh nghiệp, Gompertz Fuzzy Ranking giúp kết hợp linh hoạt đầu ra của nhiều mô hình và tăng độ ổn định khi doanh nghiệp nằm gần ranh giới giữa các nhóm xếp hạng.
\subsubsection{Dynamic Model Fusion} Dynamic Model Fusion (DMF) là phương pháp kết hợp mô hình theo hướng động, trong đó kết quả cuối cùng được điều chỉnh theo mức độ tin cậy của từng mô hình đối với từng mẫu dữ liệu. Khác với các phương pháp voting hoặc averaging sử dụng trọng số cố định, DMF kết hợp với Dynamic Classifier Selection (DCS) để lựa chọn mô hình phù hợp hơn trong từng trường hợp \cite{toufiq_explainable_2026}. Trong nghiên cứu này, DMF được sử dụng để kết hợp hai mô hình T-LSTM và GraphSAGE. T-LSTM có ưu thế trong khai thác xu hướng tài chính theo thời gian, trong khi GraphSAGE hỗ trợ học quan hệ đồ thị giữa các doanh nghiệp. Gọi $\hat{y}_{T}$ và $\hat{y}_{G}$ lần lượt là nhãn dự đoán của T-LSTM và GraphSAGE. Nếu hai mô hình cho cùng một nhãn, kết quả cuối cùng được xác định như sau: 
\begin{equation} 
\hat{y} = \hat{y}_{T} = \hat{y}_{G} 
\end{equation} Ngược lại, nếu hai mô hình dự đoán khác nhau, hệ thống tính điểm tin cậy động cho từng mô hình dựa trên độ tin cậy trên tập validation và xác suất dự đoán của mẫu kiểm thử: 
\begin{equation} 
R_m(x) = \lambda V_m(\hat{y}_m) + (1-\lambda) P_m(\hat{y}_m|x), 
\end{equation} trong đó $R_m(x)$ là điểm tin cậy của mô hình $m$, $V_m(\hat{y}_m)$ là độ tin cậy của mô hình trên tập validation đối với lớp được dự đoán, $P_m(\hat{y}_m|x)$ là xác suất dự đoán của mô hình đối với mẫu $x$, và $\lambda$ là hệ số cân bằng. Nhãn cuối cùng được chọn từ mô hình có điểm tin cậy cao hơn: 
\begin{equation} 
\hat{y} = \hat{y}_{m^*}, \quad m^* = \arg\max_{m \in \{T,G\}} R_m(x). 
\end{equation} Nhờ cơ chế lựa chọn động, DMF giúp tận dụng thế mạnh bổ sung giữa mô hình chuỗi thời gian và mô hình đồ thị, từ đó nâng cao độ ổn định trong bài toán xếp hạng tín dụng doanh nghiệp.
\subsection{Các phương pháp Explainable AI - XAI}
\subsubsection{GradientSHAP} 
GradientSHAP là phương pháp giải thích dựa trên gradient, kết hợp ý tưởng của SHAP với cách lấy mẫu ngẫu nhiên giữa dữ liệu cần giải thích và các mẫu nền. Phương pháp này ước lượng mức đóng góp của từng đặc trưng bằng cách tính gradient của đầu ra mô hình theo đầu vào tại nhiều điểm nội suy giữa mẫu cần giải thích và các mẫu baseline \cite{lundberg_unified_2017, kokhlikyan_captum_2020}. 
\subsubsection{LIME} 
LIME là phương pháp giải thích cục bộ độc lập với mô hình, được Ribeiro và cộng sự đề xuất nhằm giải thích một dự báo cụ thể bằng một mô hình đơn giản tại vùng lân cận của mẫu cần phân tích \cite{ribeiro_why_2016}. Ý tưởng chính của LIME là tạo ra nhiều mẫu nhiễu xung quanh quan sát gốc, đưa các mẫu này vào mô hình để lấy xác suất dự báo, sau đó huấn luyện một mô hình tuyến tính cục bộ có trọng số để xấp xỉ hành vi của mô hình gốc.
\section{PHƯƠNG PHÁP ĐỀ XUẤT}
Để giải quyết bài toán, tôi đề xuất sử dụng mô hình tổng quát như hình dưới đây. Mô hình tổng quát gồm có 2 pha: Pha huấn luyện và pha kiểm thử. được mô tả cụ  
thể như Hình 1.
\begin{figure}[htbp]
    \centering % Căn giữa hình ảnh
    \includegraphics[width=1\textwidth]{uploads/sodo.drawio.drawio.png} % Thay 'ten-file-anh' bằng tên ảnh của bạn (không cần đuôi .png/.jpg)
    \caption{Mô hình đề xuất} % Dòng chú thích dưới ảnh
\end{figure}
\subsection{Pha huấn luyện} 
\subsubsection{Tiền xử lý dữ liệu và chuẩn bị dữ liệu} 
Trong pha huấn luyện, dữ liệu tài chính doanh nghiệp được tiền xử lý nhằm đảm bảo chất lượng đầu vào, tăng độ ổn định và cải thiện khả năng hội tụ của các mô hình. Trước hết, các bản ghi trùng lặp theo mã chứng khoán và ngày xếp hạng được gộp lại, trong đó các chỉ số tài chính được lấy theo giá trị trung vị, còn nhãn xếp hạng được chọn theo mức rủi ro cao hơn nhằm đảm bảo nguyên tắc thận trọng. Các giá trị khuyết thiếu ở biến số được điền bằng trung vị của tập huấn luyện để hạn chế ảnh hưởng của ngoại lai và tránh rò rỉ dữ liệu. Tiếp theo, các đặc trưng tài chính có phân phối lệch được xử lý bằng phép biến đổi logarit khi độ lệch vượt ngưỡng cho trước, giúp dữ liệu gần với phân phối chuẩn hơn. Sau đó, toàn bộ 12 chỉ số tài chính được chuẩn hóa bằng phương pháp Z-score nhằm đưa các biến về cùng hệ quy chiếu, giảm sự chi phối của những đặc trưng có thang đo lớn và hỗ trợ quá trình tối ưu hóa của mô hình. Do bộ dữ liệu ban đầu gồm 21 hạng tín dụng có sự mất cân bằng lớn giữa các lớp, các hạng này được gộp thành ba nhóm chính gồm Investment Grade (IG), High Yield (HY) và Distressed. Cách gộp này giúp giảm độ thưa của nhãn, đồng thời vẫn giữ được ý nghĩa phân biệt rủi ro tín dụng. Cuối cùng, dữ liệu được sắp xếp theo từng doanh nghiệp và thời gian để tạo các cửa sổ chuỗi thời gian có độ dài cố định. Đối với các lớp thiểu số, TimeGAN được sử dụng để sinh thêm chuỗi tài chính tổng hợp, góp phần cân bằng tập huấn luyện và nâng cao khả năng tổng quát hóa của hệ thống.
\subsubsection{Rút trích đặc trưng và huấn luyện}
Trong các kịch bản thực nghiệm, quá trình rút trích đặc trưng chuỗi thời gian được thực hiện tự động và trực tiếp thông qua các mô hình học sâu dựa trên cơ chế cửa sổ trượt. Cấu trúc đầu vào bao gồm các chỉ số tài chính đã qua chuẩn hóa, đặc trưng gia số và các biến phân loại được nhúng. Quá trình huấn luyện sử dụng thuật toán tối ưu AdamW kết hợp với hàm mất mát tùy chỉnh Cross Entropy được thiết kế kết hợp Cross Entropy đặc thù nhằm khắc phục tình trạng mất cân bằng dữ liệu và tối ưu hóa tính thứ bậc của các nhãn xếp hạng tín nhiệm. Để kiểm soát quá trình hội tụ và ngăn ngừa hiện tượng quá khớp, kỹ thuật Gradient Clipping được áp dụng, điều chỉnh dừng huấn luyện dựa trên sự cải thiện của chỉ số F1. Tất cả các mô hình được huấn luyện độc lập với cùng một quy trình tiền xử lý dữ liệu thống nhất, đảm bảo tính nhất quán khi so sánh. Trong suốt quá trình huấn luyện, các chỉ số như Loss, Accuracy được theo dõi và ghi nhận liên tục, phục vụ cho việc đánh giá, đối chiếu và báo cáo hiệu suất mô hình.
\subsection{Pha kiểm thử} Sau khi hoàn tất quá trình huấn luyện, các mô hình được đưa vào pha kiểm thử nhằm đánh giá khả năng tổng quát hóa trên dữ liệu chưa từng thấy. Trong giai đoạn này, mô hình được chuyển sang chế độ đánh giá bằng hàm \texttt{eval()} để vô hiệu hóa các thành phần như Dropout và Batch Normalization, từ đó đảm bảo kết quả dự đoán ổn định và phản ánh đúng hiệu năng thực tế. Đầu vào của pha kiểm thử là các chuỗi dữ liệu tài chính đã được tiền xử lý, chuẩn hóa và cấu trúc hóa theo từng doanh nghiệp. Mô hình thực hiện dự đoán ba nhóm xếp hạng tín dụng gồm Investment Grade (IG), High Yield (HY) và Distressed. Kết quả dự đoán được đánh giá thông qua các chỉ số tiêu chuẩn như Accuracy, Loss, Precision, Recall, F1-Score, AUC-ROC và ma trận nhầm lẫn. Bên cạnh các mô hình đơn lẻ, nghiên cứu còn kiểm thử nhiều kịch bản kết hợp mô hình nhằm khai thác thế mạnh bổ sung giữa các kiến trúc khác nhau. Các tổ hợp Fuzzy Ensemble bao gồm TTX (TCN, T-LSTM, XGBoost), PLL (PatchTST, LSTM, LightGBM) và TTLPXL (T-LSTM, TCN, LSTM, PatchTST, XGBoost, LightGBM). Các tổ hợp này được triển khai theo hai phương pháp kết hợp là Fuzzy Choquet Integral và Gompertz Fuzzy Ranking nhằm tổng hợp xác suất dự đoán của các mô hình thành phần một cách linh hoạt hơn so với voting hoặc averaging truyền thống. Ngoài ra, nghiên cứu xây dựng kịch bản Dynamic Model Fusion (DMF) kết hợp T-LSTM và GraphSAGE. Trong đó, T-LSTM khai thác xu hướng biến động tài chính theo thời gian, còn GraphSAGE học các quan hệ cấu trúc trong dữ liệu doanh nghiệp. Khi hai mô hình cho kết quả khác nhau, cơ chế Dynamic Classifier Selection được sử dụng để lựa chọn mô hình có độ tin cậy cao hơn dựa trên hiệu năng validation và mức độ tự tin của dự đoán. Để nâng cao tính minh bạch, pha kiểm thử tích hợp các kỹ thuật Explainable AI gồm GradientSHAP và LIME. GradientSHAP được sử dụng để đánh giá mức độ đóng góp của từng chỉ số tài chính đối với kết quả dự báo, trong khi LIME giải thích cục bộ từng trường hợp doanh nghiệp cụ thể. Nhờ đó, hệ thống không chỉ đưa ra nhãn xếp hạng cuối cùng mà còn cung cấp cơ sở diễn giải cho quyết định phân loại, góp phần nâng cao độ tin cậy và khả năng ứng dụng thực tiễn của mô hình.


\section{KẾT QUẢ THỰC NGHIỆM}

\subsection{Môi trường cài đặt và tập dữ liệu thực nghiệm}

\subsubsection{Môi trường cài đặt}
Các thử nghiệm trong nghiên cứu được thực hiện trên nền tảng Kaggle Notebooks, hỗ trợ môi trường tính toán đám mây phù hợp cho các bài toán máy học và học sâu \cite{noauthor_notebooks_nodate}. Hệ thống sử dụng CPU Intel Xeon 2.00 GHz với 4 vCPU, 32 GB RAM và GPU NVIDIA Tesla P100 với 16 GB VRAM nhằm tăng tốc quá trình huấn luyện. Các mô hình được triển khai bằng Python 3.8+ cùng các thư viện mã nguồn mở phổ biến. Trong đó, PyTorch được sử dụng làm framework chính cho các mô hình học sâu nhờ khả năng tính toán linh hoạt và hỗ trợ GPU hiệu quả. Pandas và NumPy được dùng để xử lý dữ liệu, trong khi Matplotlib, Scikit-learn và Scikit-metrics hỗ trợ trực quan hóa và đánh giá mô hình. Môi trường Kaggle giúp thuận tiện trong triển khai và tái lập kết quả, tuy nhiên giới hạn thời gian sử dụng GPU yêu cầu quá trình huấn luyện được tối ưu về batch size và bộ nhớ.

\subsubsection{Tập dữ liệu thực nghiệm} Tập dữ liệu thực nghiệm được tổng hợp từ hai bộ dữ liệu công khai trên Kaggle, gồm \textit{Corporate Credit Rating} \cite{noauthor_corporate_credit_rating_nodate} và \textit{Corporate Credit Rating with Financial Ratios} \cite{noauthor_corporate_financial_ratios_nodate}. Hai bộ dữ liệu được tích hợp thông qua các trường khóa như tên công ty, mã chứng khoán, tổ chức xếp hạng và ngày xếp hạng. Sau quá trình làm sạch, gộp dữ liệu và loại bỏ các bản ghi trùng lặp, tập dữ liệu cuối cùng gồm 8.680 bản ghi quan sát thuộc 1.623 doanh nghiệp trong giai đoạn 2005--2016. Mỗi bản ghi được biểu diễn bởi 12 chỉ số tài chính gồm Current Ratio, Debt-to-Equity Ratio, Gross Profit Margin, Operating Profit Margin , EBIT Margin, Pretax Profit Margin, Net Profit Margin, ROE, ROA, Asset Turnover, Operating, Cashflow per Share, Free Cashflow per Share. Phản ánh các khía cạnh quan trọng về sức khỏe tài chính của doanh nghiệp như thanh khoản, đòn bẩy, khả năng sinh lời, hiệu quả hoạt động và dòng tiền. Các chỉ số này được tính toán từ báo cáo tài chính, bao gồm bảng cân đối kế toán, báo cáo kết quả kinh doanh và báo cáo lưu chuyển tiền tệ. Do các hạng tín dụng ban đầu có sự mất cân bằng lớn, 21 hạng xếp hạng được gộp thành ba nhóm chính gồm Investment Grade (IG), High Yield (HY) và Distressed. Tập dữ liệu có đặc điểm là dữ liệu bảng kết hợp yếu tố chuỗi thời gian, trong đó mỗi doanh nghiệp có trung bình khoảng 5,35 lần được xếp hạng. Đặc điểm này phù hợp với các mô hình học sâu xử lý chuỗi như LSTM, TCN, PatchTST và T-LSTM. Để hạn chế rò rỉ dữ liệu, quá trình chia tập huấn luyện, validation và kiểm thử được thực hiện theo mã chứng khoán độc lập. Thông tin chi tiết về số lượng mẫu của từng lớp sau khi gộp nhãn được trình bày trong Bảng~\ref{tab:class_distribution}. 
\begin{table}[H] 
\centering 
\caption{Thông tin chi tiết về số dữ liệu của từng lớp ban đầu} 
\label{tab:class_distribution} 
\begin{tabular}{lccc} 
\hline \textbf{Lớp} & \textbf{Train} & \textbf{Val} & \textbf{Test} \\ \hline Investment Grade (IG) & 3865 & 552 & 1105 \\ High Yield (HY) & 1957 & 281 & 282 \\ Distressed & 207 & 30 & 29 \\ \hline 
\end{tabular} 
\end{table}
\subsection{Kịch bản thực nghiệm}
Tổng cộng có 14 kịch bản thực nghiệm được triển khai, được chia thành hai nhóm chính. Nhóm đầu tiên gồm 7 kịch bản mô hình đơn, mô hình lai và mô hình đồ thị (kịch bản từ 1 đến 7), bao gồm: LightGBM, XGBoost, Temporal Convolutional Networks (TCN), Long Short-Term Memory (LSTM), Patch Time Series Transformer (PatchTST), Transformer-Long Short-Term Memory (T-LSTM), GraphSAGE. Nhóm thứ hai gồm 7 kịch bản kết hợp mô hình (kịch bản từ 8 đến 14), trong đó các mô hình được kết hợp theo 3 phương pháp học kết hợp: Tích phân mờ Choquet (Fuzzy Choquet Integral – FI), Xếp hạng mờ Gompertz (Gompertz Fuzzy Ranking – FR) và Dynamic Model Fusion (DMF) được trình bày trong Bảng~\ref{tab:experimental_scenarios}.
\begin{table}[H] 
\centering 
\caption{Các kịch bản thực nghiệm} 
\label{tab:experimental_scenarios} 
\begin{tabular}{c l l >{\raggedright\arraybackslash}p{6.5cm}} 
\hline 
\textbf{Kịch bản} & \textbf{Tên kịch bản} & \textbf{Phương pháp kết hợp} & \textbf{Mô hình mạng sử dụng} \\ \hline 
1 & LightGBM & -- & LightGBM \\ 
2 & XGBoost & -- & XGBoost \\ 
3 & TCN & -- & Temporal Convolutional Networks (TCN) \\ 
4 & LSTM & -- & Long Short-Term Memory (LSTM) \\ 
5 & PatchTST & -- & Patch Time Series Transformer (PatchTST) \\ 
6 & T-LSTM & -- & Transformer--Long Short-Term Memory (T-LSTM) \\ 
7 & GraphSAGE & -- & GraphSAGE, Multilayer Perceptron (MLP) \\ 
8 & FI-TTX & Fuzzy Choquet Integral & TCN + T-LSTM + XGBoost \\ 
9 & FI-PLL & Fuzzy Choquet Integral & PatchTST + LSTM + LightGBM \\ 
10 & FI-TTLPXL & Fuzzy Choquet Integral & TCN + T-LSTM + LSTM + PatchTST + XGBoost + LightGBM \\ 
11 & FR-TTX & Gompertz Fuzzy Ranking & TCN + T-LSTM + XGBoost \\ 
12 & FR-PLL & Gompertz Fuzzy Ranking & PatchTST + LSTM + LightGBM \\ 
13 & FR-TTLPXL & Gompertz Fuzzy Ranking & TCN + T-LSTM + LSTM + PatchTST + XGBoost + LightGBM \\ 
14 & DMF & Dynamic Model Fusion & GraphSAGE + T-LSTM \\ \hline \end{tabular} 
\end{table}
Với các kịch bản 7 kịch bản mô hình đơn, mô hình lai và mô hình đồ thị (kịch bản từ 1 đến 7) được các tham số huấn luyện sau.
\begin{table}[H] 
\centering 
\caption{Tham số huấn luyện các kịch bản} 
\label{tab:hyperparameter} 
\begin{tabular}{c c c l c} 
\hline 
\textbf{Kịch bản} & \textbf{Learning rate (Tốc độ học)} & \textbf{Epochs (Số lần học)} & \textbf{Loss (Hàm mất mát)} & \textbf{Num classes (Số lớp)} \\ 
\hline 
1 & 3e-2 & 100 & Cross Entropy & 12 \\ 
2 & 3e-2 & 100 & Cross Entropy & 12 \\ 
3 & 3e-2 & 100 & Cross Entropy & 12 \\ 
4 & 3e-4 & 100 & Cross Entropy & 12 \\ 
5 & 3e-4 & 100 & Cross Entropy & 12 \\ 
6 & 3e-4 & 100 & Cross Entropy & 12 \\ 
7 & 3e-2 & 100 & Cross Entropy & 12 \\ 
\hline 
\end{tabular} 
\end{table}
\subsection{Kết quả huấn luyện}

\begin{table}[H]
\centering
\caption{Bảng tổng hợp kết quả huấn luyện của các mô hình (\%)}
\label{tab:training_results}
\begin{tabular}{c l c c c c}
\hline
\textbf{Kịch bản} & \textbf{Mạng huấn luyện} & \textbf{Train Loss} & \textbf{Val Loss} & \textbf{Train Acc} & \textbf{Val Acc} \\
\hline
1 & LightGBM  & 0.123 & 0.266 & 94.0 & 92.0 \\
2 & XGBoost   & 0.177 & 0.267 & 94.0 & 92.4 \\
3 & TCN       & 0.248 & 0.276 & 92.7 & 92.3 \\
4 & LSTM      & 0.190 & 0.268 & 94.0 & 92.2 \\
5 & PatchTST  & 0.217 & 0.279 & 93.2 & 91.7 \\
6 & T-LSTM    & 0.167 & 0.223 & 94.0 & 92.4 \\
7 & GraphSAGE & 0.343 & 0.399 & 92.1 & 92.2 \\
\hline
\end{tabular}
\end{table}
\subsubsection{Độ chính xác – Accuracy}
\begin{figure}[htbp]
    \centering % Căn giữa hình ảnh
    \includegraphics[width=1\textwidth]{uploads/train_val_accuracy_comparison.png}  
    \caption{Biểu đồ độ chính xác} % Dòng chú thích dưới ảnh
    \label{tab:acc-train}
\end{figure}
Theo Hình~\ref{tab:acc-train} cho thấy độ chính xác của các mô hình trên tập huấn luyện dao động từ 92,10\% đến 94,00\%, trong khi trên tập xác thực dao động từ 91,70\% đến 92,40\%. Trong đó, LightGBM, XGBoost, LSTM và T-LSTM đạt độ chính xác huấn luyện cao nhất với 94,00\%. Trên tập xác thực, XGBoost và T-LSTM cho kết quả tốt nhất, cùng đạt 92,40\%. TCN đạt 92,70\% trên tập huấn luyện và 92,30\% trên tập xác thực, cho thấy hiệu suất ổn định giữa hai tập dữ liệu. GraphSAGE cũng duy trì kết quả cân bằng với 92,10\% ở tập huấn luyện và 92,20\% ở tập xác thực. PatchTST đạt 93,20\% trên tập huấn luyện và 91,70\% trên tập xác thực, là mô hình có độ chính xác xác thực thấp nhất. Nhìn chung, các mô hình có kết quả khá đồng đều, với khoảng biến thiên trên tập xác thực chỉ 0,70 điểm phần trăm. Xét riêng theo độ chính xác, T-LSTM là mô hình nổi bật do đạt hiệu suất cao và ổn định.

\begin{table}[H]
\centering
\caption{Độ chính xác các kịch bản}
\label{tab:accuracy_models}

\renewcommand{\arraystretch}{1.05}
\setlength{\tabcolsep}{0pt}

\begin{tabular}{>{\centering\arraybackslash}p{0.48\textwidth}|
                >{\centering\arraybackslash}p{0.48\textwidth}}
\hline

\textbf{Kịch bản 1: LightGBM} &
\textbf{Kịch bản 2: XGBoost} \\
\includegraphics[width=0.85\linewidth]{acc-loss/lighgbm.png} &
\includegraphics[width=0.85\linewidth]{acc-loss/acc-xgboost.png} \\
\hline

\textbf{Kịch bản 3: TCN} &
\textbf{Kịch bản 4: LSTM} \\
\includegraphics[width=0.85\linewidth]{acc-loss/acc-tcn.png} &
\includegraphics[width=0.85\linewidth]{acc-loss/acc-lstm.png} \\
\hline

\textbf{Kịch bản 5: PatchTST} &
\textbf{Kịch bản 6: T-LSTM} \\
\includegraphics[width=0.85\linewidth]{acc-loss/acc-patchtst.png} &
\includegraphics[width=0.85\linewidth]{acc-loss/acc-tlstm.png} \\
\hline

\textbf{Kịch bản 7: GraphSAGE} &
 \\
\includegraphics[width=0.85\linewidth]{acc-loss/acc-graph.png} &
 \\
\hline

\end{tabular}
\end{table}

\subsubsection{Độ mất mát – Loss}
\begin{figure}[htbp]
    \centering % Căn giữa hình ảnh
    \includegraphics[width=1\textwidth]{acc-loss/train_val_loss_comparison.png} 
    \caption{Biểu đồ độ mất mát} % Dòng chú thích dưới ảnh
\end{figure}

Kết quả đánh giá cho thấy loss trên tập huấn luyện dao động từ 0,123 đến 0,343, trong khi loss trên tập xác thực dao động từ 0,223 đến 0,399. Ở tất cả các mô hình, validation loss đều cao hơn train loss, phản ánh xu hướng mô hình học tốt hơn trên dữ liệu huấn luyện so với dữ liệu chưa thấy. LightGBM đạt train loss thấp nhất với 0,123, tuy nhiên validation loss tăng lên 0,266, cho thấy vẫn tồn tại chênh lệch giữa khả năng học và khả năng tổng quát hóa. T-LSTM là mô hình có validation loss thấp nhất, đạt 0,223, đồng thời train loss cũng thấp ở mức 0,167. Kết quả này cho thấy T-LSTM có khả năng học ổn định và tổng quát hóa tốt trên dữ liệu tài chính dạng chuỗi thời gian. Các mô hình XGBoost, LSTM, TCN và PatchTST có validation loss tương đối gần nhau, lần lượt là 0,267, 0,268, 0,276 và 0,279. GraphSAGE có loss cao nhất trên cả tập huấn luyện và tập xác thực, lần lượt là 0,343 và 0,399, cho thấy xác suất dự báo của mô hình chưa thật sự sắc nét. Nhìn chung, xét theo độ mất mát, T-LSTM là mô hình nổi bật nhất nhờ đạt validation loss thấp nhất và duy trì hiệu suất ổn định.

\begin{table}[H]
\centering
\caption{Độ mất mát các kịch bản}
\label{tab:loss_models}

\renewcommand{\arraystretch}{1.05}
\setlength{\tabcolsep}{0pt}

\begin{tabular}{>{\centering\arraybackslash}p{0.48\textwidth}|
                >{\centering\arraybackslash}p{0.48\textwidth}}
\hline

\textbf{Kịch bản 1: LightGBM} &
\textbf{Kịch bản 2: XGBoost} \\
\includegraphics[width=0.75\linewidth]{acc-loss/loss-lightgbm.png} &
\includegraphics[width=0.75\linewidth]{acc-loss/loss-xgboost.png} \\
\hline

\textbf{Kịch bản 3: TCN} &
\textbf{Kịch bản 4: LSTM} \\
\includegraphics[width=0.75\linewidth]{acc-loss/loss-tcn.png} &
\includegraphics[width=0.75\linewidth]{acc-loss/loss-lstm.png} \\
\hline

\textbf{Kịch bản 5: PatchTST} &
\textbf{Kịch bản 6: T-LSTM} \\
\includegraphics[width=0.75\linewidth]{acc-loss/loss-oatchtst.png} &
\includegraphics[width=0.75\linewidth]{acc-loss/loss-tlstm.png} \\
\hline

\textbf{Kịch bản 7: GraphSAGE} &
 \\
\includegraphics[width=0.75\linewidth]{acc-loss/loss-graph.png} &
 \\
\hline

\end{tabular}
\end{table}

\subsubsection{Thời gian huấn luyện}

\begin{figure}[htbp]
    \centering % Căn giữa hình ảnh
    \includegraphics[width=1\textwidth]{acc-loss/time-train.png} % Thay 'ten-file-anh' bằng tên ảnh của bạn (không cần đuôi .png/.jpg)
    \caption{Thời gian huấn luyện các kịch bản } % Dòng chú thích dưới ảnh
\end{figure}
Kết quả thời gian huấn luyện cho thấy sự khác biệt rõ rệt giữa các mô hình, phản ánh mức độ phức tạp và chi phí tính toán của từng kiến trúc. T-LSTM là mô hình có thời gian huấn luyện cao nhất với 7,07 phút, tiếp theo là TCN với 5,83 phút, do các mô hình này phải khai thác quan hệ phụ thuộc theo thời gian trong dữ liệu tài chính. PatchTST, LSTM và GraphSAGE có thời gian huấn luyện ở mức trung bình, lần lượt là 3,42 phút, 2,42 phút và 1,70 phút. Trong khi đó, XGBoost và LightGBM huấn luyện nhanh nhất, chỉ mất 0,82 phút và 1,08 phút. Nhìn chung, các mô hình học sâu có chi phí huấn luyện cao hơn nhưng phù hợp để khai thác dữ liệu chuỗi, còn các mô hình cây tăng cường có ưu thế về tốc độ và hiệu quả tính toán.

\subsection{Kết quả kiểm thử}

\subsubsection{Độ chính xác}
\begin{figure}[htbp]
    \centering % Căn giữa hình ảnh
    \includegraphics[width=1\textwidth]{acc-loss/test_accuracy_scenario_comparison.png} 
    \caption{Biểu đồ độ chính xác trên tập kiểm thử} 
    \label{tab:acc-test}
\end{figure}
Kết quả kiểm thử cho thấy tất cả các mô hình và kịch bản kết hợp đều đạt độ chính xác trên 91\%, chứng tỏ khả năng phân loại ổn định trong bài toán xếp hạng tín dụng doanh nghiệp. Đối với các mô hình đơn lẻ, T-LSTM đạt kết quả cao nhất với 92,51\%, tiếp theo là GraphSAGE với 92,34\% và LSTM với 92,20\%. Các mô hình XGBoost, LightGBM, TCN và PatchTST đạt độ chính xác lần lượt là 92,05\%, 91,99\%, 91,93\% và 91,35\%. Trong nhóm Fuzzy Choquet Integral, FI-TTX đạt kết quả tốt nhất với 92,51\%, trong khi các kịch bản Gompertz Fuzzy Ranking có độ chính xác khá đồng đều, cao nhất là FR-TTLPXL với 92,46\%. Đáng chú ý, kịch bản Dynamic Model Fusion (DMF) đạt độ chính xác cao nhất trên toàn bộ tập kiểm thử với 92,98\%. Kết quả này cho thấy DMF tận dụng hiệu quả thế mạnh bổ sung giữa T-LSTM và GraphSAGE thông qua cơ chế lựa chọn mô hình động theo từng mẫu dữ liệu.

\subsubsection{Precision của các kịch bản}

\begin{figure}[H]
\centering


\captionsetup[subfigure]{
    labelformat=simple,
    labelsep=period,
    font=it,
    justification=centering
}

\renewcommand{\thesubfigure}{\alph{subfigure}}

\begin{subfigure}{0.72\textwidth}
    \centering
    \includegraphics[width=\textwidth]{acc-loss/distressed_precision_test_comparison.png}
    \caption{Precision của lớp Distressed}
    \label{fig:precision_distressed}
\end{subfigure}

\vspace{0.15cm}

\begin{subfigure}{0.72\textwidth}
    \centering
    \includegraphics[width=\textwidth]{acc-loss/hy_precision_test_comparison.png}
    \caption{Precision của lớp HY (High Yield)}
    \label{fig:precision_hy}
\end{subfigure}

\vspace{0.15cm}

\begin{subfigure}{0.72\textwidth}
    \centering
    \includegraphics[width=\textwidth]{acc-loss/ig_precision_test_comparison.png}
    \caption{Precision của lớp IG (Investment Grade)}
    \label{fig:precision_ig}
\end{subfigure}
\label{fig:precision_classes}

\caption{Precision lớp Distressed (a), lớp HY (b) và lớp IG (c) trên tập kiểm thử }
\end{figure}
Kết quả Precision trên tập kiểm thử cho thấy khả năng dự đoán đúng của các mô hình có sự khác biệt rõ rệt giữa từng lớp xếp hạng tín dụng. Đối với lớp Distressed, DMF đạt Precision cao nhất với 72,41\%, vượt qua các mô hình đơn lẻ và các kịch bản ensemble khác. Trong nhóm mô hình đơn, LSTM đạt kết quả tốt nhất ở lớp này với 69,23\%, trong khi FR-TTLPXL là kịch bản Fuzzy Ensemble nổi bật với 70,73\%. Đối với lớp HY, Precision dao động trong khoảng 84,41\%--87,96\%, trong đó GraphSAGE đạt cao nhất với 87,96\%, tiếp theo là DMF với 87,78\% và T-LSTM với 87,67\%. Ở lớp IG, hầu hết các mô hình đều đạt Precision rất cao và ổn định, dao động khoảng 95,81\%--96,44\%, cho thấy đây là lớp dễ nhận diện nhất. FR-TTX đạt Precision cao nhất ở lớp IG với 96,44\%. Nhìn chung, DMF cho thấy hiệu quả nổi bật ở lớp rủi ro cao Distressed, GraphSAGE và DMF thể hiện tốt ở lớp HY, trong khi các mô hình đều duy trì khả năng phân loại ổn định đối với lớp IG.

\subsubsection{Recall của các kịch bản}
\begin{figure}[H]
\centering


\captionsetup[subfigure]{
    labelformat=simple,
    labelsep=period,
    font=it,
    justification=centering
}

\renewcommand{\thesubfigure}{\alph{subfigure}}

\begin{subfigure}{0.72\textwidth}
    \centering
    \includegraphics[width=\textwidth]{acc-loss/distressed_recall_test_comparison.png}
    \caption{Recall của lớp Distressed}
    \label{fig:Recall_distressed}
\end{subfigure}

\vspace{0.15cm}

\begin{subfigure}{0.72\textwidth}
    \centering
    \includegraphics[width=\textwidth]{acc-loss/hy_recall_test_comparison.png}
    \caption{Recall của lớp HY (High Yield)}
    \label{fig:Recall_hy}
\end{subfigure}

\vspace{0.15cm}

\begin{subfigure}{0.72\textwidth}
    \centering
    \includegraphics[width=\textwidth]{acc-loss/ig_recall_test_comparison.png}
    \caption{Recall của lớp IG (Investment Grade)}
    \label{fig:Recall_ig}
\end{subfigure}
\label{fig:Recall_classes}

\caption{Recall lớp Distressed (a), lớp HY (b) và lớp IG (c) trên tập kiểm thử }
\end{figure}
Kết quả Recall cho thấy khả năng nhận diện đầy đủ các mẫu thuộc từng lớp có sự khác biệt rõ rệt giữa ba nhóm xếp hạng tín dụng. Đối với lớp Distressed, GraphSAGE đạt Recall cao nhất với 62,71\%, cho thấy mô hình đồ thị có khả năng phát hiện tốt hơn các doanh nghiệp rủi ro cao. LightGBM và XGBoost đạt lần lượt 50,85\% và 49,15\%, trong khi PatchTST có Recall thấp nhất với 20,34\%. Trong nhóm ensemble, FI-PLL, FR-PLL và FR-TTLPXL cùng đạt 49,15\%, còn DMF đạt 35,59\%. Ở lớp HY, các mô hình có Recall ổn định hơn, chủ yếu dao động quanh 89\%--91\%; trong đó DMF đạt cao nhất với 91,23\%, tiếp theo là PatchTST và FR-TTX cùng đạt 91,06\%. Đối với lớp IG, Recall của tất cả mô hình đều ở mức cao, từ 95,29\% đến 96,92\%, trong đó DMF đạt kết quả tốt nhất với 96,92\%. Nhìn chung, lớp IG là lớp dễ nhận diện nhất, lớp Distressed là lớp khó phát hiện nhất, GraphSAGE nổi bật ở khả năng nhận diện nhóm rủi ro cao, còn DMF thể hiện ưu thế ở hai lớp HY và IG.

\subsubsection{F1-Score của các kịch bản}
\begin{figure}[H]
\centering


\captionsetup[subfigure]{
    labelformat=simple,
    labelsep=period,
    font=it,
    justification=centering
}

\renewcommand{\thesubfigure}{\alph{subfigure}}

\begin{subfigure}{0.72\textwidth}
    \centering
    \includegraphics[width=\textwidth]{acc-loss/distressed_f1_score_test_comparison.png}
    \caption{F1-Score của lớp Distressed}
    \label{fig:F1-Score_distressed}
\end{subfigure}

\vspace{0.15cm}

\begin{subfigure}{0.72\textwidth}
    \centering
    \includegraphics[width=\textwidth]{acc-loss/hy_f1_score_test_comparison.png}
    \caption{F1-Score của lớp HY (High Yield)}
    \label{fig:F1-Score_hy}
\end{subfigure}

\vspace{0.15cm}

\begin{subfigure}{0.72\textwidth}
    \centering
    \includegraphics[width=\textwidth]{acc-loss/ig_f1_score_test_comparison.png}
    \caption{F1-Score của lớp IG (Investment Grade)}
    \label{fig:F1-Score_ig}
\end{subfigure}
\label{fig:F1-Score_classes}

\caption{F1-Score lớp Distressed (a), lớp HY (b) và lớp IG (c) trên tập kiểm thử }
\end{figure}
F1-score là chỉ số kết hợp giữa Precision và Recall, phản ánh mức độ cân bằng giữa khả năng dự đoán đúng và khả năng nhận diện đầy đủ các mẫu thuộc từng lớp. Kết quả thực nghiệm cho thấy lớp Distressed có sự khác biệt lớn nhất giữa các mô hình, trong đó GraphSAGE đạt F1-score cao nhất với 63,25\%, thể hiện khả năng cân bằng tốt khi nhận diện nhóm doanh nghiệp rủi ro cao. Trong nhóm ensemble, FR-TTLPXL đạt kết quả tốt nhất ở lớp Distressed với 58,00\%, trong khi PatchTST có F1-score thấp nhất với 30,00\%. Đối với lớp HY, F1-score dao động trong khoảng 87,61\%--89,47\%, cho thấy các mô hình duy trì hiệu quả phân loại ổn định; DMF đạt kết quả cao nhất với 89,47\%, tiếp theo là T-LSTM với 88,99\%. Ở lớp IG, tất cả mô hình đều đạt F1-score rất cao, từ 95,55\% đến 96,57\%, trong đó DMF tiếp tục dẫn đầu với 96,57\%. Nhìn chung, DMF nổi bật ở hai lớp HY và IG, còn GraphSAGE thể hiện tốt nhất ở lớp Distressed. Điều này cho thấy việc đánh giá mô hình xếp hạng tín dụng cần xem xét hiệu quả theo từng lớp rủi ro, thay vì chỉ dựa vào độ chính xác tổng thể.

\subsubsection{Ma trận nhầm lẫn}

\begin{table}[H]
\centering
\caption{Ma trận nhầm lẫn các kịch bản trên tập kiểm thử}
\label{tab:confusion_matrices}

\renewcommand{\arraystretch}{1.2}
\setlength{\tabcolsep}{4pt}

\begin{tabular}{|>{\centering\arraybackslash}m{0.31\textwidth}|
                >{\centering\arraybackslash}m{0.31\textwidth}|
                >{\centering\arraybackslash}m{0.31\textwidth}|}
\hline

\includegraphics[width=0.95\linewidth]{acc-loss/cm-lightgbm.png} &
\includegraphics[width=0.95\linewidth]{acc-loss/cm-xgboost.png} &
\includegraphics[width=0.95\linewidth]{acc-loss/cm-tcn.png} \\
\hline

\includegraphics[width=0.95\linewidth]{acc-loss/cm-lstm.png} &
\includegraphics[width=0.95\linewidth]{acc-loss/cm-patchtst.png} &
\includegraphics[width=0.95\linewidth]{acc-loss/cm-tlstm.png} \\
\hline

\includegraphics[width=0.95\linewidth]{acc-loss/cm-graph.png} & 
\includegraphics[width=0.95\linewidth]{acc-loss/cm-fittx.png} &
\includegraphics[width=0.95\linewidth]{acc-loss/cm-fipll.png} \\
\hline

\includegraphics[width=0.95\linewidth]{acc-loss/cm-fittlpxl.png} & 
\includegraphics[width=0.95\linewidth]{acc-loss/cm-frttx.png} &
\includegraphics[width=0.95\linewidth]{acc-loss/cm-frpll.png} \\
\hline

\includegraphics[width=0.95\linewidth]{acc-loss/cm-ttlpxl.png} & & \\
\hline

\end{tabular}
\end{table}

Kết quả ma trận nhầm lẫn cho thấy các mô hình phân loại tốt hai lớp HY và IG, trong đó lớp IG được nhận diện ổn định nhất với phần lớn số mẫu nằm trên đường chéo chính. Ngược lại, lớp Distressed là lớp khó phân loại nhất, chủ yếu bị nhầm sang lớp HY do ranh giới đặc trưng giữa hai nhóm rủi ro này chưa thật sự rõ ràng. Trong các mô hình đơn lẻ, GraphSAGE dự đoán đúng nhiều mẫu Distressed nhất với 37 mẫu, cho thấy ưu thế trong phát hiện nhóm rủi ro cao. T-LSTM đạt hiệu quả tốt ở lớp IG với 1068 mẫu dự đoán đúng. Đáng chú ý, DMF/DCS đạt hiệu quả tổng thể cao nhất khi dự đoán đúng 510 mẫu HY và 1071 mẫu IG, tuy nhiên khả năng nhận diện lớp Distressed vẫn thấp hơn GraphSAGE.

\subsubsection{AUC-ROC của các kịch bản}


\begin{figure}[H]
\centering


\captionsetup[subfigure]{
    labelformat=simple,
    labelsep=period,
    font=it,
    justification=centering
}

\renewcommand{\thesubfigure}{\alph{subfigure}}

\begin{subfigure}{0.72\textwidth}
    \centering
    \includegraphics[width=\textwidth]{acc-loss/combined_auc_roc_distressed_14_models.png}
    \caption{Đường cong ROC trên lớp Distressed}
    \label{fig:AUC-ROC_distressed}
\end{subfigure}

\vspace{0.15cm}

\begin{subfigure}{0.72\textwidth}
    \centering
    \includegraphics[width=\textwidth]{acc-loss/combined_auc_roc_hy_14_models.png}
    \caption{Đường cong ROC trên lớp HY (High Yield)}
    \label{fig:AUC-ROC_hy}
\end{subfigure}

\vspace{0.15cm}

\begin{subfigure}{0.72\textwidth}
    \centering
    \includegraphics[width=\textwidth]{acc-loss/combined_auc_roc_ig_14_models.png}
    \caption{Đường cong ROC trên lớp IG (Investment Grade)}
    \label{fig:AUC-ROC_ig}
\end{subfigure}
\label{fig:AUC-ROC_classes}

\caption{AUC-ROC lớp Distressed (a), lớp HY(b) và lớp IG (c) trên tập kiểm thử}
\end{figure}
Kết quả AUC-ROC cho thấy hầu hết các mô hình đều có khả năng phân biệt tốt giữa từng lớp xếp hạng và các lớp còn lại, thể hiện qua các đường cong ROC nằm xa đường chéo ngẫu nhiên. Đối với lớp Distressed, FR-TTX đạt AUC cao nhất với 96,1\%, tiếp theo là XGBoost và FR-TTLPXL cùng đạt 96,0\%. Ở lớp HY, FI-TTX cho kết quả tốt nhất với AUC 97,4\%, cho thấy hiệu quả của Fuzzy Ensemble trong phân biệt nhóm rủi ro trung bình. Đối với lớp IG, T-LSTM đạt AUC cao nhất với 98,8\%, phản ánh khả năng nhận diện rất tốt nhóm doanh nghiệp có chất lượng tín dụng cao. Nhìn chung, các kịch bản Fuzzy Ensemble như FI-TTX, FR-TTX và FR-TTLPXL có AUC nổi bật ở nhiều lớp, trong khi DMF đạt độ chính xác tổng thể cao nhưng chưa phải là mô hình có AUC tốt nhất theo từng lớp.

\subsubsection{Thời gian kiểm thử}
\begin{figure}[htbp]
    \centering % Căn giữa hình ảnh
    \includegraphics[width=1\textwidth]{acc-loss/test_inference_time_comparison.png} % Thay 'ten-file-anh' bằng tên ảnh của bạn (không cần đuôi .png/.jpg)
    \caption{Thời gian kiểm thử của các các kịch bản } % Dòng chú thích dưới ảnh
\end{figure}

Kết quả thời gian kiểm thử cho thấy tất cả mô hình và kịch bản kết hợp đều có tốc độ suy luận rất nhanh, dao động từ 0,101 giây đến 0,169 giây. Trong đó, FI-PLL và DMF đạt thời gian kiểm thử thấp nhất, cùng 0,101 giây, tiếp theo là FI-TTX với 0,104 giây và FR-PLL với 0,109 giây. Các mô hình đơn lẻ như T-LSTM, GraphSAGE và LSTM cũng duy trì thời gian kiểm thử thấp, lần lượt đạt 0,121 giây, 0,129 giây và 0,136 giây. Ngược lại, PatchTST và FR-TTX có thời gian kiểm thử cao nhất, cùng đạt 0,169 giây, trong khi FR-TTLPXL đạt 0,160 giây. Nhìn chung, mức chênh lệch giữa các kịch bản là không lớn, cho thấy toàn bộ hệ thống đáp ứng tốt yêu cầu về tốc độ suy luận và có tiềm năng triển khai trong thực tế.

\subsubsection{Explainable AI của các kịch bản}
\paragraph{GradientSHAP}
\begin{table}[H]
\centering
\caption{Ma trận nhầm lẫn các kịch bản trên tập kiểm thử}
\label{tab:confusion_matrices}

\renewcommand{\arraystretch}{1.2}
\setlength{\tabcolsep}{4pt}

\begin{tabular}{|>{\centering\arraybackslash}m{0.31\textwidth}|
                >{\centering\arraybackslash}m{0.31\textwidth}|
                >{\centering\arraybackslash}m{0.31\textwidth}|}
\hline

\includegraphics[width=0.95\linewidth]{acc-loss/cm-lightgbm.png} &
\includegraphics[width=0.95\linewidth]{acc-loss/cm-xgboost.png} &
\includegraphics[width=0.95\linewidth]{acc-loss/cm-tcn.png} \\
\hline

\includegraphics[width=0.95\linewidth]{acc-loss/cm-lstm.png} &
\includegraphics[width=0.95\linewidth]{acc-loss/cm-patchtst.png} &
\includegraphics[width=0.95\linewidth]{acc-loss/cm-tlstm.png} \\
\hline

\includegraphics[width=0.95\linewidth]{acc-loss/cm-graph.png} & 
\includegraphics[width=0.95\linewidth]{acc-loss/cm-fittx.png} &
\includegraphics[width=0.95\linewidth]{acc-loss/cm-fipll.png} \\
\hline

\includegraphics[width=0.95\linewidth]{acc-loss/cm-fittlpxl.png} & 
\includegraphics[width=0.95\linewidth]{acc-loss/cm-frttx.png} &
\includegraphics[width=0.95\linewidth]{acc-loss/cm-frpll.png} \\
\hline

\includegraphics[width=0.95\linewidth]{acc-loss/cm-ttlpxl.png} & & \\
\hline

\end{tabular}
\end{table}
\section{KẾT LUẬN}


\bibliographystyle{ieeetr}
\bibliography{thesis}

\end{document}
